We propose Process-Aware Policy Optimization (\method{}), a method that integrates process-level evaluation into Group Relative Policy Optimization (GRPO) through \textbf{decoupled advantage normalization}, to address two limitations of existing reward designs. Outcome reward models (ORM) evaluate only final-answer correctness, treating all correct responses identically regardless of reasoning quality, and gradually lose the advantage signal as groups become uniformly correct. Process reward models (PRM) offer richer supervision, but directly using PRM scores causes reward hacking, where models exploit verbosity to inflate scores while accuracy collapses. \method{} resolves both by composing the advantage from an outcome component $\aout$, derived from ORM and normalized over all responses, and a process component $\aproc$, derived from a rubric-based PRM and normalized exclusively among correct responses. This decoupled design ensures that $\aout$ anchors training on correctness while $\aproc$ differentiates reasoning quality without distorting the outcome signal. Experiments across multiple model scales and six benchmarks demonstrate that \method{} consistently outperforms ORM, reaching 51.3\% vs.\ 46.3\% on OlympiadBench while continuing to improve as ORM plateaus and declines.
